\label{sec:dense_method}
We fine-tune an OpenVLA-OFT-based policy with SimpleVLA-RL and LoRA on complete LIBERO-Spatial instructions~\cite{openvlaoft,simplevlarl,libero}. The reward comparison separates terminal-only feedback, phase-completion feedback, and a mixture of phase-completion bonuses with continuous within-phase progress. Thus the dense condition already includes intermediate phase rewards; phase-only shaping removes the continuous progress component.

A tracker follows reach, grasp, lift, move, and place. Continuous progress uses gripper--object distance, gripper closure, object height, object--target distance, and release state. Within the active phase, progress credit is given only when it exceeds the previous maximum, which resets at each phase transition. Progress increments and phase-completion events are weighted and combined, and the shaping contribution is bounded before combination with terminal feedback. The native task-success verifier remains unchanged.

The combined dense-and-phase condition reaches 65.3\% success versus 61.7\% for terminal-only reward at the reported 160-update budget (Table~\ref{tab:libero}). Phase-only shaping reaches 62.3\%, showing that completed subgoals can also provide useful feedback. Section~\ref{sec:libero_results} distinguishes these comparisons and their settings.

For full-task composition, we choose milestone-based feedback to reduce task-specific reward design across long-horizon instructions. Extending dense shaping requires a progress measure and its normalization for each additional action, together with the relevant object, target, and contact conditions. Milestones require completion predicates but avoid defining continuous progress throughout every phase. We therefore retain dense rewards for acquiring individual skills and use binary milestones to compose the complete task. This choice targets a more reusable reward structure; the reported results do not establish that milestone-only feedback outperforms every dense configuration.
